\documentclass{article}
\usepackage{amsmath}  % For mathematical symbols
\usepackage{amssymb}  % For set notation like \mathbb

\begin{document}

\title{Hierarchy of Numbering Systems}
\author{}
\date{}
\maketitle

\section*{Hierarchy of Numbering Systems}

\begin{itemize}
    \item \textbf{Natural Numbers} (\(\mathbb{N}\)): 
    \[
    \mathbb{N} = \{1, 2, 3, 4, \dots \}
    \]
    (Some definitions also include 0: \(\{0, 1, 2, 3, \dots\}\)).
    
    \item \textbf{Whole Numbers}: 
    \[
    \{0, 1, 2, 3, \dots \}
    \]
    
    \item \textbf{Integers} (\(\mathbb{Z}\)): 
    \[
    \mathbb{Z} = \{\dots, -3, -2, -1, 0, 1, 2, 3, \dots \}
    \]
    
    \item \textbf{Rational Numbers} (\(\mathbb{Q}\)): 
    \[
    \mathbb{Q} = \left\{\frac{a}{b} : a, b \in \mathbb{Z}, b \neq 0 \right\}
    \]
    Rational numbers include all numbers that can be expressed as a ratio of two integers.
    
    \item \textbf{Irrational Numbers}: 
    \[
    \{ \sqrt{2}, \pi, e, \dots \}
    \]
    Irrational numbers are numbers that cannot be expressed as a ratio of two integers, with non-repeating and non-terminating decimal expansions.
    
    \item \textbf{Real Numbers} (\(\mathbb{R}\)): 
    \[
    \mathbb{R} = \text{Rational Numbers} \cup \text{Irrational Numbers}
    \]
    The real numbers include both rational and irrational numbers and represent every point on the continuous number line.
    
    \item \textbf{Imaginary Numbers} (\(\mathbb{I}\)): 
    \[
    \mathbb{I} = \{ i \cdot b : i = \sqrt{-1}, b \in \mathbb{R} \}
    \]
    Imaginary numbers involve the square root of negative numbers, where \(i = \sqrt{-1}\).
    
    \item \textbf{Complex Numbers} (\(\mathbb{C}\)): 
    \[
    \mathbb{C} = \{ a + bi : a, b \in \mathbb{R}, i = \sqrt{-1} \}
    \]
    Complex numbers combine real and imaginary numbers in the form \(a + bi\).
    
\end{itemize}

\end{document}
